\question
\tocsection{25 Horses}

There are 25 horses of unknown speed. They can be raced in groups of five to find the finishing order, but their times are not recorded. What is the minimum number of races necessary to determine the fastest three horses and what strategy should be used to accomplish this?

\begin{hints}
	\begin{enumerate}
		\item Can you think of an easier problem involving the 25 horses?
		\item How would you determine the fastest horse?
		\item The second-fastest horse may have been in the same race as the fastest horse. Keep track of where the possible fast horses may be and then race them together.
		\item Seven races are necessary and sufficient.
	\end{enumerate}
\end{hints}

\begin{solution}

All horses need to be raced at some point and 25 naturally decomposes into five races of five horses which we will call the heats. The fastest horse must be among the winners of the heats and racing these five together in a sixth race determines the winner. We label the heats A, B, C, D, and E in order of how fast the fastest horse in each race was and we label the horses within each race by their position. For example, B3 is the horse that came third in the heat with the horse that came second in the sixth race.

The second fastest horse would have either won their heat or have been beaten by the fastest horse if they shared the same heat. This means the second fastest horse was either A2 or B1. We cannot distinguish between these cases with the information available and there is clearly no other way of finding the fastest horse in six races. We therefore conclude that a seventh race is necessary. The third fastest horse could be any of the following three:

\begin{enumerate}
	\item A3: This is where the second fastest horse was in the same heat as the fastest horse and both of them beat the third fastest.
	\item B2: This is where the second fastest horse was in a different heat from the fastest horse, but shared a heat with the third fastest.
	\item C1: The top three fastest horses were in different heats so all made it through to the sixth race.
\end{enumerate}

Together this makes A2, B1, A3, B1, and C1 which is five horses. If these horses are raced together then the top two positions give the second and third fastest horses. The different possibilities for the final race are given in table~\ref{tab:25Horses_FinalRace}.

\begin{table}[H]
	\centering
	\caption{Final race possibilities based on the initial position of the fastest three horses.}
	\label{tab:25Horses_FinalRace}
	\begin{tabular}{ccc@{\hspace{8mm}}cc}
		\toprule
		Race 1 & Race 2 & Race 3 & Second Horse & Third Horse  \\
		\midrule
		1, 2, 3 &      &   & A2 & A3  \\
		1, 2    & 3    &   & A2 & B1  \\
		1       & 2, 3 &   & B1 & B2  \\
		1, 3    & 2    &   & B1 & A2  \\
		1       & 2    & 3 & B1 & C1  \\
		\bottomrule
	\end{tabular}
\end{table}

\vspace{-4ex}
\phantom{ }

\end{solution}